\subsection{Hypothesis 6 Validation}
The results in the statistical analysis of hypothesis 6 demonstrates an effect of cultural group across the difference between emotional and analytical trust. Because the WEIRD/Normal participant grouping is determined by multiple participant demographics, there is the potential of a hidden mediating effect of one of these demographics individually. If this were the case, the main conclusions of hypothesis 6 would not be supported and instead the observed effect would be better described as differences related to one or more of the constituent demographics. To test the possibility of this, we performed mediation analyses testing whether there was a mediating effect of each of the five demographic items on the impact of the demographic group on analytical vs. emotional trust shown in Figure \ref{fig:hypothesis6}. The results below demonstrate that none of the constituent demographics display a mediating effect for either of the two conditions, providing further support to hypothesis 6. 

Across both conditions combined (Interactive and Text), separate mediation models tested whether each WEIRD-defining attribute mediated the association between WEIRD group membership and the analytical-versus-emotional trust differences. For each model, $X$ denotes the WEIRD indicator, $M$ denotes the candidate mediator, and $Y=\Delta A-\Delta E$; education was coded as above bachelor versus bachelor or below and ethnicity as white, other, or prefer not to say versus black, asian, or mixed; $b$ denotes the path $(M \to Y\mid X)$, $c$ denotes the path $X\rightarrow Y$, and $c'$ denotes the path $(X \to Y\mid M)$. Indirect effects were estimated using non-parametric bootstrap implemented in Python. 

For the \textbf{western geographic indicator} (residence or origin) as mediator, analysis with $n=454$ showed an $a$ path of $=0.6819$, a $b$ path of $=-0.0250$, a total effect of $c=-0.0770$, and a direct effect of $c'=-0.0599$. The estimated indirect effect was $ab=-0.0168$, with $95\%$ CI $[-0.1019,\,0.0672]$ and $p=.686$. This indicates no clear evidence of mediation because the bootstrap confidence interval includes zero. 

For the \textbf{english language indicator} as mediator, analysis with $n=454$ showed an $a$ path of $=0.4971$, a $b$ path of $=-0.1024$, a total effect of $c=-0.0770$, and a direct effect of $c'=-0.0261$. The estimated indirect effect was $ab=-0.0509$, with $95\%$ CI $[-0.1091,\,0.0050]$ and $p=.073$. This indicates no clear evidence of mediation because the bootstrap confidence interval includes zero. 

For the \textbf{education indicator} as mediator, analysis with $n=454$ showed an $a$ path of $=0.1119$, a $b$ path of $=-0.0655$, a total effect of $c=-0.0770$, and a direct effect of $c'=-0.0697$. The estimated indirect effect was $ab=-0.0072$, with $95\%$ CI $[-0.0279,\,0.0059]$ and $p=.336$. This indicates no clear evidence of mediation because the bootstrap confidence interval includes zero.

For the \textbf{employment indicator} (employed full-time or not in paid work vs.\ all other employment statuses grouped together) as mediator, analysis with $n=454$ showed an $a$ path of $=0.4286$, a $b$ path of $=-0.0088$, a total effect of $c=-0.0770$, and a direct effect of $c'=-0.0732$. The estimated indirect effect was $ab=-0.0042$, with $95\%$ CI $[-0.0512,\,0.0408]$ and $p=.862$. This indicates no clear evidence of mediation because the bootstrap confidence interval includes zero. 

For the \textbf{ethnicity indicator} as mediator, analysis with $n=454$ showed an $a$ path of $=0.3140$, a $b$ path of $=0.0987$, a total effect of $c=-0.0770$, and a direct effect of $c'=-0.1080$. The estimated indirect effect was $ab=0.0308$, with $95\%$ CI $[-0.0030,\,0.0687]$ and $p=.070$. This indicates no clear evidence of mediation because the bootstrap confidence interval includes zero.

Overall these five analyses demonstrate that the impact that cultural background has on the difference between emotional and analytical trust is not better explained as a mediated effect of any of the constituent demographics used in our categorization of WEIRD and Normal participants.

Table \ref{tab:mediation_all} reports the same mediation analysis for every pairwise comparison of the groups within each demographic variable (country of origin excluded).

{\footnotesize
\setlength{\tabcolsep}{2pt}
\newcolumntype{R}[1]{>{\raggedleft\arraybackslash}p{#1}}
\newcolumntype{L}[1]{>{\raggedright\arraybackslash}p{#1}}
\begin{longtable}{L{0.12\linewidth}L{0.25\linewidth}R{0.05\linewidth}R{0.065\linewidth}R{0.065\linewidth}R{0.065\linewidth}R{0.16\linewidth}R{0.055\linewidth}R{0.055\linewidth}}
\caption{Mediation of the effect of WEIRD group on the analytical-versus-emotional trust-change difference for every pairwise comparison of each demographic variable (country of origin excluded). Only participants in the two compared groups are included; $M=1$ for the first group listed and $M=0$ for the second (reference) group. $a$: $X\rightarrow M$; $b$: $M\rightarrow Y\mid X$; $ab$: indirect effect with 95\% bootstrap CI (5000 resamples); $p$: bootstrap $p$-value; $q$: Benjamini-Hochberg adjusted $p$ within each demographic variable. Bold rows have a bootstrap CI that excludes zero.}
\label{tab:mediation_all}\\
\toprule
Variable & Comparison & $n$ & $a$ & $b$ & $ab$ & 95\% CI & $p$ & $q$ \\
\midrule
\endfirsthead
\toprule
Variable & Comparison & $n$ & $a$ & $b$ & $ab$ & 95\% CI & $p$ & $q$ \\
\midrule
\endhead
Age group & 18-24 vs. 25-34 & 308 & -0.288 & -0.042 & 0.012 & [-0.026, 0.056] & .550 & .688 \\
Age group & 35-44 vs. 25-34 & 283 & 0.083 & -0.060 & -0.005 & [-0.028, 0.011] & .573 & .688 \\
Age group & 45+ vs. 25-34 & 261 & 0.219 & -0.159 & -0.035 & [-0.089, 0.003] & .073 & .439 \\
Age group & 35-44 vs. 18-24 & 193 & 0.455 & -0.023 & -0.010 & [-0.088, 0.071] & .802 & .802 \\
Age group & 45+ vs. 18-24 & 171 & 0.591 & -0.115 & -0.067 & [-0.176, 0.040] & .210 & .631 \\
Age group & 45+ vs. 35-44 & 146 & 0.173 & -0.099 & -0.017 & [-0.064, 0.015] & .336 & .673 \\
\midrule
Gender & Man vs. Woman & 187 & -0.197 & 0.024 & -0.004 & [-0.044, 0.029] & .787 & .944 \\
Gender & Non-binary vs. Woman & 165 & -0.083 & 0.197 & -0.017 & [-0.080, 0.029] & .470 & .940 \\
Gender & Rather not say vs. Woman & 158 & 0.044 & -0.019 & -0.001 & [-0.034, 0.027] & .946 & .946 \\
Gender & Non-binary vs. Man & 156 & 0.119 & 0.191 & 0.024 & [-0.037, 0.104] & .433 & .940 \\
Gender & Rather not say vs. Man & 149 & 0.249 & -0.039 & -0.009 & [-0.083, 0.052] & .754 & .944 \\
Gender & Rather not say vs. Non-binary & 127 & 0.134 & -0.206 & -0.029 & [-0.120, 0.030] & .352 & .940 \\
\midrule
AI knowledge & Advanced vs. Conceptual understanding & 329 & -0.033 & -0.037 & 0.001 & [-0.010, 0.015] & .860 & .950 \\
AI knowledge & Beginner knowledge vs. Conceptual understanding & 327 & -0.000 & 0.028 & -0.000 & [-0.011, 0.010] & .950 & .950 \\
AI knowledge & Expert vs. Conceptual understanding & 230 & 0.073 & -0.277 & -0.023 & [-0.101, 0.016] & .368 & .950 \\
AI knowledge & Beginner knowledge vs. Advanced & 220 & 0.037 & 0.068 & 0.002 & [-0.015, 0.025] & .815 & .950 \\
AI knowledge & Expert vs. Advanced & 123 & 0.147 & -0.187 & -0.030 & [-0.148, 0.050] & .506 & .950 \\
AI knowledge & Expert vs. Beginner knowledge & 121 & 0.127 & -0.253 & -0.034 & [-0.154, 0.035] & .392 & .950 \\
\midrule
\textbf{Education} & \textbf{High School vs. Bachelor} & \textbf{315} & \textbf{-0.275} & \textbf{0.166} & \textbf{-0.046} & \textbf{[-0.092, -0.008]} & \textbf{.016} & \textbf{.160} \\
Education & Master vs. Bachelor & 303 & -0.003 & -0.027 & -0.000 & [-0.011, 0.010] & .981 & .981 \\
Education & Graduate Professional Degree vs. Bachelor & 256 & -0.046 & 0.002 & -0.001 & [-0.019, 0.015] & .908 & .981 \\
Education & PhD vs. Bachelor & 231 & 0.135 & -0.055 & -0.008 & [-0.078, 0.046] & .837 & .981 \\
\textbf{Education} & \textbf{Master vs. High School} & \textbf{184} & \textbf{0.398} & \textbf{-0.166} & \textbf{-0.067} & \textbf{[-0.152, -0.002]} & \textbf{.039} & \textbf{.196} \\
Education & Graduate Professional Degree vs. High School & 137 & 0.340 & -0.126 & -0.044 & [-0.164, 0.039] & .329 & .981 \\
Education & PhD vs. High School & 112 & 0.642 & -0.059 & -0.039 & [-0.220, 0.152] & .658 & .981 \\
Education & Graduate Professional Degree vs. Master & 125 & -0.077 & 0.020 & -0.002 & [-0.040, 0.030] & .884 & .981 \\
Education & PhD vs. Master & 100 & 0.265 & 0.029 & 0.007 & [-0.120, 0.129] & .886 & .981 \\
Education & PhD vs. Graduate Professional Degree & 53 & 0.468 & 0.069 & 0.032 & [-0.167, 0.261] & .751 & .981 \\
\midrule
Employment status & Part-Time vs. Full-Time & 281 & -0.302 & -0.101 & 0.029 & [-0.011, 0.073] & .151 & .378 \\
Employment status & Unemployed (and job seeking) vs. Full-Time & 274 & -0.387 & 0.086 & -0.034 & [-0.095, 0.026] & .260 & .443 \\
Employment status & Other vs. Full-Time & 216 & -0.122 & 0.175 & -0.021 & [-0.054, 0.004] & .099 & .356 \\
Employment status & Not in paid work vs. Full-Time & 202 & 0.025 & -0.140 & -0.004 & [-0.027, 0.012] & .731 & .812 \\
Employment status & Unemployed (and job seeking) vs. Part-Time & 181 & -0.372 & 0.174 & -0.066 & [-0.155, 0.003] & .065 & .356 \\
Employment status & Other vs. Part-Time & 123 & 0.070 & 0.273 & 0.020 & [-0.062, 0.133] & .684 & .812 \\
Employment status & Not in paid work vs. Part-Time & 109 & 0.368 & -0.013 & -0.005 & [-0.121, 0.117] & .907 & .907 \\
Employment status & Other vs. Unemployed (and job seeking) & 116 & 0.518 & 0.094 & 0.045 & [-0.087, 0.198] & .469 & .669 \\
Employment status & Not in paid work vs. Unemployed (and job seeking) & 102 & 0.762 & -0.178 & -0.146 & [-0.452, 0.115] & .266 & .443 \\
Employment status & Not in paid work vs. Other & 44 & 0.410 & -0.321 & -0.131 & [-0.358, 0.019] & .107 & .356 \\
\midrule
\textbf{Ethnicity} & \textbf{Black vs. White} & \textbf{263} & \textbf{-0.206} & \textbf{-0.163} & \textbf{0.035} & \textbf{[0.004, 0.078]} & \textbf{.021} & \textbf{.198} \\
\textbf{Ethnicity} & \textbf{Asian vs. White} & \textbf{262} & \textbf{-0.253} & \textbf{-0.176} & \textbf{0.045} & \textbf{[0.004, 0.095]} & \textbf{.026} & \textbf{.198} \\
Ethnicity & Mixed vs. White & 243 & -0.278 & -0.054 & 0.015 & [-0.031, 0.060] & .503 & .980 \\
Ethnicity & Other vs. White & 233 & -0.178 & -0.128 & 0.024 & [-0.014, 0.076] & .270 & .810 \\
Ethnicity & Prefer not to say vs. White & 205 & -0.028 & -0.081 & 0.001 & [-0.020, 0.022] & .944 & .980 \\
Ethnicity & Asian vs. Black & 149 & -0.107 & 0.002 & -0.000 & [-0.034, 0.030] & .980 & .980 \\
Ethnicity & Mixed vs. Black & 130 & -0.350 & 0.142 & -0.053 & [-0.138, 0.006] & .095 & .474 \\
Ethnicity & Other vs. Black & 120 & -0.075 & 0.045 & -0.003 & [-0.045, 0.038] & .871 & .980 \\
Ethnicity & Prefer not to say vs. Black & 92 & 0.143 & 0.083 & 0.008 & [-0.079, 0.108] & .854 & .980 \\
Ethnicity & Mixed vs. Asian & 129 & -0.300 & 0.130 & -0.043 & [-0.131, 0.014] & .181 & .680 \\
Ethnicity & Other vs. Asian & 119 & 0.024 & 0.036 & 0.002 & [-0.043, 0.050] & .936 & .980 \\
Ethnicity & Prefer not to say vs. Asian & 91 & 0.240 & 0.113 & 0.025 & [-0.091, 0.175] & .656 & .980 \\
Ethnicity & Other vs. Mixed & 100 & 0.368 & -0.117 & -0.050 & [-0.203, 0.047] & .409 & .980 \\
Ethnicity & Prefer not to say vs. Mixed & 72 & 0.606 & 0.090 & 0.047 & [-0.198, 0.296] & .647 & .980 \\
Ethnicity & Prefer not to say vs. Other & 62 & 0.259 & 0.037 & 0.014 & [-0.128, 0.193] & .870 & .980 \\
\midrule
Language & English vs. Non-English & 454 & 0.497 & -0.102 & -0.051 & [-0.109, 0.005] & .073 & .073 \\
\bottomrule
\end{longtable}
}


\subsubsection{Fine-grained Analysis}
\textcolor{red}{To preform a more fine-grained analysis of the impact of inclusion criteria on the difference web observed between participants based on cultural group, we performed an extensive mediation analysis. This fine-grained mediation analysis compared every variable in the inclusion criteria for WEIRD vs. Normal to determine if any sub-group of one of these demographics had a mediation effect. Out of 54 different mediation analyses (full table in supplementary materials) the only significant mediation effects observed were in two sub groupings of Education and Ethnicity). However, these mediation effects did not remain significant after B-H correction, as shown in the following q values. These mediation effects within education were for High School vs. Bachelor ($n=315, a=-0.275, b=0.166, ab = -0.046, 95\%[-0.092, -0.008], p =0.039, q=0.160$) and for Master's vs. High School ($n=184, a=0.398, b=-0.166, ab=-0.067, 95\%[-0.152, -0.002], p =0.016$ with a q-value of $p=0.196$). These mediation effects within ethnicity were for Black vs. White ($n=263, a=-0.206, b=-0.163, ab=0.035, 95\%[0.004, 0.078], p=0.021$) with a q-value of $0.198$ and for Asian vs. White ($n=262, a=-0.253, b=-0.176, ab=0.045,95\%[-0.092, -0.008], p=0.026$) with a q-value of $0.198$. See the supplementary materials for the full information on all mediation tests run. }
